\documentclass[portrait,color=sotonbracinggreen,margin=3cm]{sotonposter}

% \usepackage{svg} % uncomment if you include SVG figures

\title{Poster title goes here, \\ \vspace{3mm} split across two lines}

\author[1 *]{First Author}
\author[1]{Second Author}
\affil[1]{Research Group, Campus, University of Southampton. \protect\\}
\affil[*]{first.author@soton.ac.uk}

\begin{document}

\maketitle
\vspace{-42mm}

\begin{multicols}{2}

    \section*{Section title}

    % Coloured highlight boxes are the main building block of this
    % template. Pass any colour as the optional argument; mixing with
    % white (e.g. sotonbryellow!30!white) softens it for backgrounds.
    \begin{highlightbox}[sotonbryellow!30!white]
        \textit{\textbf{Background}}:
        Introduce the project here. Two to four sentences work best:
        what the problem is, why it matters, and what this poster
        shows. Keep it readable from a distance.
    \end{highlightbox}

    \begin{minipage}{\linewidth}
        \centering
        \includegraphics[width=0.9\linewidth]{example-image-a}
        \captionof{figure}{A schematic or overview figure introducing
        the problem. Replace \texttt{example-image-a} with your own
        figure from the shared figures folder.}
        \label{fig:schematic}
    \end{minipage}

    \columnbreak

    \begin{highlightbox}[white!10!white]
        \textit{\textbf{Methods}}: describe the approach here, citing
        tools or prior work as needed\cite{ExampleReference}. Refer to
        figures in the usual way, e.g. Figure~\ref{fig:schematic}.
    \end{highlightbox}

    \vspace{-3mm}

    Plain text also works between boxes, for details that do not need
    highlighting: parameter values, test conditions, assumptions and
    similar supporting information.

    \begin{highlightbox}[sotonneonblue!30!white]
        \textit{\textbf{Equations}} sit naturally inside boxes:
        \begin{equation}
            f(x) = a_0 \, \exp\!\left( -\frac{(x - x_0)^2}{2\sigma^2}
            \right) \sin( 2 \pi f t )
        \end{equation}
        with symbol definitions given directly below the equation.
    \end{highlightbox}

    \begin{highlightbox}[sotonbgreen!10!white]
        \textit{\textbf{Key concept}}: use a box like this to explain
        the central idea of the work, with a pointer to the main \cite{FayRiddell1958}
        result, e.g. Figure~\ref{fig:mainresult}.
    \end{highlightbox}

    \vspace{-8mm}

\end{multicols}

% Full-width figure spanning both columns, for the main result.
\vspace{-10mm}
\begin{minipage}{\linewidth}
    \centering
    \captionsetup{justification=centering,margin=2cm}
    \includegraphics[width=0.6\linewidth]{example-image}
    \captionof{figure}{A wide, full-width figure for the main result
    of the poster.}
    \label{fig:mainresult}
\end{minipage}

\vspace{-1mm}

\begin{multicols}{2}

    \begin{highlightbox}[sotonlilac!80!white]
        \textit{\textbf{Analysis}}: a second row of columns for
        results and discussion, for example
        \begin{equation}
            N = \ln\!\left( \frac{A(x)}{A(x_0)} \right)
        \end{equation}
        with the terms defined in the surrounding text.
    \end{highlightbox}



    \columnbreak



    \begin{highlightbox}[sotonbpurple!20!white]
        \textit{\textbf{Conclusions}}: summarise the findings and any
        acknowledgements of funding or support here.
    \end{highlightbox}

\end{multicols}

\vspace{-23mm}
\bibliography{references}

\end{document}